\section{Additional related works}
\label{app:addition_related_works}
\begin{table}[ht]
\centering
\caption{Positioning of H-SFP relative to resource-efficient, split, and prototype-based FL methods. H-SFP is designed for hierarchical split-FL and replaces sample-wise cross-tier communication with class-wise feature statistics.}
\label{tab:method_positioning}
\resizebox{\linewidth}{!}{
\begin{tabular}{lcccccc}
\toprule
Method & Hierarchical & Split model & Cross-tier BP & Transmitted object & Generator-free & Formal privacy \\
\midrule
FedAvg & \multirow{3}{*}{\xmark} & \multirow{3}{*}{\xmark} & \multirow{3}{*}{N/A} & \multirow{3}{*}{full model updates} & \multirow{3}{*}{\cmark} & \multirow{3}{*}{\xmark} \\
FedProx & & & & & & \\
FedNova & & & & & & \\
\midrule
SplitFed / HSFL & \cmark & \cmark & \cmark & activations + gradients & \cmark & \xmark \\
FedProto & \xmark & \xmark & N/A & class prototypes & \cmark & \xmark \\
FedGen / FedDF & \xmark & \xmark & N/A & generator/logit information & \xmark & \xmark \\
PriSM & partial & \cmark & method-specific & sampled kernels/sub-models & \cmark & \xmark \\
Delta & partial & \cmark & reduced/asymmetric & private/residual flows & \cmark & \cmark \\
\textbf{H-SFP (Ours)} & \cmark & \cmark & \xmark & class-wise $(\mu,\sigma)$ & \cmark & \xmark \\
\bottomrule
\end{tabular}}
\end{table}